\documentclass[a4paper,12pt]{article} % Dokumenttyp: 'article' für kurze Arbeiten

% ----- Pakete -----
\usepackage[utf8]{inputenc}    % Zeichencodierung (Umlaute etc.)
\usepackage[T1]{fontenc}       % Verbesserte Zeichendarstellung
\usepackage[ngerman]{babel}    % Deutsche Spracheinstellungen
\usepackage{amsmath, amssymb}  % Mathematische Symbole und Umgebungen
\usepackage{graphicx}          % Zum Einbinden von Bildern
\usepackage{geometry}          % Seitenränder anpassen
\usepackage{hyperref}          % Klickbare Links und PDF-Metadaten
\usepackage{enumitem}          % Bessere Kontrolle über Aufzählungen
\usepackage{booktabs}          % Schönere Tabellen
\usepackage{xcolor}            % Farben für Text, Tabellen etc.

% ----- Layout -----
\geometry{a4paper, margin=2.5cm}  % Seitenränder

% ----- PDF-Metadaten -----
\hypersetup{
    colorlinks=true,
    linkcolor=blue,
    urlcolor=teal,
    pdftitle={Projektbericht Simulation einer Wasserrakete},
    pdfauthor={Simon Kerlin, Patrick Stempel},
}

% ----- Titel -----
\title{Projektbericht\\ Simulation einer Wasserrakete}
\author{Simon Kerlin, Patrick Stempel}
% \date{\today}

% ==========================
\begin{document}
% ==========================

\maketitle % Titel ausgeben

\tableofcontents % Inhaltsverzeichnis (optional)
\newpage

\section{Einleitung}
Hier beginnt dein Text. Du kannst normale Absätze einfach durch eine Leerzeile trennen.

\section{Mathematische Umgebung}
Ein Beispiel für eine Gleichung:
\begin{equation}
E = mc^2
\end{equation}

Oder inline wie \( a^2 + b^2 = c^2 \).

\section{Aufzählungen}
\begin{itemize}
    \item Punkt 1
    \item Punkt 2
\end{itemize}

Oder nummeriert:
\begin{enumerate}
    \item Erster Punkt
    \item Zweiter Punkt
\end{enumerate}

\section{Bilder}
\begin{figure}[h!]
    \centering
    % \includegraphics[width=0.5\textwidth]{beispielbild.png}
    \caption{Ein Beispielbild.}
    \label{fig:beispiel}
\end{figure}

\section{Tabellen}
\begin{table}[h!]
    \centering
    \begin{tabular}{l c r}
        \toprule
        Links & Mitte & Rechts \\
        \midrule
        A & B & C \\
        D & E & F \\
        \bottomrule
    \end{tabular}
    \caption{Beispieltabelle}
    \label{tab:beispiel}
\end{table}

\section{Fazit}
Hier schließt du dein Dokument ab.

% ==========================
\end{document}
% ==========================
